
\documentclass{article}
\usepackage{colortbl}
\usepackage{makecell}
\usepackage{multirow}
\usepackage{supertabular}

\begin{document}

\newcounter{utterance}

\centering \large Interaction Transcript for game `codenames', experiment `risk\_low', episode 4 with reflexion{-}Qwen3.5{-}27B.
\vspace{24pt}

{ \footnotesize  \setcounter{utterance}{1}
\setlength{\tabcolsep}{0pt}
\begin{supertabular}{c@{$\;$}|p{.15\linewidth}@{}p{.15\linewidth}p{.15\linewidth}p{.15\linewidth}p{.15\linewidth}p{.15\linewidth}}
    \# & $\;$A & \multicolumn{4}{c}{Game Master} & $\;\:$B\\
    \hline

    \theutterance \stepcounter{utterance}  
    & & & \multicolumn{2}{p{0.3\linewidth}}{
        \cellcolor[rgb]{0.9,0.9,0.9}{
            \makecell[{{p{\linewidth}}}]{
                \texttt{\tiny{[GM$|$GM]}}
                {'hidden': {'team': ['canada', 'pitch', 'state'], 'innocent': ['mail', 'berlin', 'club', 'cap', 'web', 'train'], 'opponent': [], 'assassin': []}, 'revealed': {'team': {'team': [], 'innocent': [], 'opponent': [], 'assassin': []}, 'opponent': {'team': [], 'innocent': [], 'opponent': [], 'assassin': []}}}
            }
        }
    }
    & & \\ \\

    \theutterance \stepcounter{utterance}  
    & & \multicolumn{4}{p{0.6\linewidth}}{
        \cellcolor[rgb]{0.9,0.9,0.9}{
            \makecell[{{p{\linewidth}}}]{
                \texttt{\tiny{[P1$\langle$GM]}}
                \texttt{You have attempted this same task before and did not succeed. These are the plans you wrote after those attempts.} \\
\texttt{Trial \#1: New plan: The previous attempt failed because I used "mouse" as the clue, which is explicitly forbidden by the rule stating the clue "cannot be one of the words in the lists." Additionally, I incorrectly included "diamond" in my first turn's targets, which was an opponent word, causing an early turn end. Next time, I will strictly verify that my chosen clue is not present in the team, opponent, distractor, or kill word lists. I will also ensure that the TARGETS list only contains words from my current remaining team words, avoiding any opponent or distractor words to prevent ending the turn prematurely. For the current state, I will choose a valid clue like "mountain" for "olympus" and "canada" (or "pitch" for "baseball" context if applicable, though "mountain" is safer for olympus/canada) and list only valid team words.} \\
\\ 
\texttt{Trial \#2: New plan: My previous strategy failed because I prioritized grouping multiple words over ensuring the clue was distinct enough from the opponent's and distractor's lists. Specifically, "baseball" led to the opponent word "field," and "money" led to the distractor "mine." I also failed to account for the fact that "province" might be too abstract or ambiguous for "state" in this context, or that the turn ended before "canada" was guessed. Next time, I will strictly analyze the semantic overlap between my potential clue and the *opponent* and *distractor* lists before submitting. If a clue risks triggering an opponent word (like "baseball" {-}> "field") or a distractor (like "money" {-}> "mine"), I will either choose a more specific clue or split the targets into separate turns to ensure safety. For the remaining words "canada, pitch, state," I will avoid broad categories like "province" or "money" and instead use a clue that uniquely targets "canada" (e.g., "maple") or "state" (e.g., "federal") individually to prevent accidental guesses of opponent/distractor words, ensuring I clear the board without ending the turn prematurely.} \\
\\ 
\\ 
\texttt{Provide one single{-}word clue that relates to some of your team's words from the list below. You can choose to how many and to which words your clue relates to.} \\
\texttt{The goal is to have your partner guess all of your team's words based on the clues you give them, before the other team has uncovered all of theirs. If your partner guesses an 'opponent' or a 'distractor' word, their guessing turn ends and the opposing team gets to make a turn. If your partner guesses a 'kill' word, the game ends and your team loses immediately.} \\
\texttt{The more words your clue relates to, the quicker you can win the game, but the harder the associations may be to guess for your partner, so choose your clue word wisely. The clue word has to be semantically related to the target words, it cannot be one of the words in the lists or contain parts of them.} \\
\\ 
\texttt{Always give your single{-}word clue and your comma{-}separated list of related target words in the following format and make your answers as short as possible, never include any other text than is required in this form:} \\
\\ 
\texttt{CLUE: <WORD>} \\
\texttt{TARGETS: <TARGETS>} \\
\\ 
\texttt{Your team words are: canada, pitch, state.} \\
\texttt{Your opponent's team words are: .} \\
\texttt{Distractor words are: mail, berlin, club, cap, web, train.} \\
\texttt{The kill words are: .} \\
            }
        }
    }
    & \\ \\

    \theutterance \stepcounter{utterance}  
    & \multicolumn{4}{p{0.6\linewidth}}{
        \cellcolor[rgb]{0.9,0.9,0.9}{
            \makecell[{{p{\linewidth}}}]{
                \texttt{\tiny{[P1$\rangle$GM]}}
                \texttt{CLUE: maple} \\
\texttt{TARGETS: canada} \\
            }
        }
    }
    & & \\ \\

    \theutterance \stepcounter{utterance}  
    & & & \multicolumn{2}{p{0.3\linewidth}}{
        \cellcolor[rgb]{0.9,0.9,0.9}{
            \makecell[{{p{\linewidth}}}]{
                \texttt{\tiny{[GM$|$GM]}}
                \texttt{maple} \\
            }
        }
    }
    & & \\ \\

    \theutterance \stepcounter{utterance}  
    & & & \multicolumn{2}{p{0.3\linewidth}}{
        \cellcolor[rgb]{0.9,0.9,0.9}{
            \makecell[{{p{\linewidth}}}]{
                \texttt{\tiny{[GM$|$GM]}}
                ['canada']
            }
        }
    }
    & & \\ \\

    \theutterance \stepcounter{utterance}  
    & & & \multicolumn{2}{p{0.3\linewidth}}{
        \cellcolor[rgb]{0.9,0.9,0.9}{
            \makecell[{{p{\linewidth}}}]{
                \texttt{\tiny{[GM$|$GM]}}
                {'word': 'canada', 'assignment': 'team'}
            }
        }
    }
    & & \\ \\

    \theutterance \stepcounter{utterance}  
    & & & \multicolumn{2}{p{0.3\linewidth}}{
        \cellcolor[rgb]{0.9,0.9,0.9}{
            \makecell[{{p{\linewidth}}}]{
                \texttt{\tiny{[GM$|$GM]}}
                \texttt{opponent has won} \\
            }
        }
    }
    & & \\ \\

    \theutterance \stepcounter{utterance}  
    & \multicolumn{4}{p{0.6\linewidth}}{
        \cellcolor[rgb]{0.95,0.95,0.95}{
            \makecell[{{p{\linewidth}}}]{
                \texttt{\tiny{[P2$\rightarrow$P1]}}
                \texttt{Post{-}episode Reflexion (not sent to GM):} \\
\texttt{New plan: My previous attempt failed because I only targeted one word ("canada") with the clue "maple," leaving "pitch" and "state" unguessed, which resulted in a loss as the game requires clearing all team words. While "maple" was a safe clue, the strategy of isolating words was too slow. Next time, I will prioritize finding a single clue that semantically links multiple remaining team words to maximize efficiency. For the words "canada," "pitch," and "state," I will look for a broader concept like "nation" (linking Canada and State) or "sports" (linking Pitch and potentially State if interpreted as a team, though risky). If a multi{-}word link is too ambiguous and risks a distractor, I will still aim to group at least two words (e.g., "maple" for Canada and "federal" for State) rather than just one, ensuring I make significant progress toward clearing the board in fewer turns. I will strictly verify that the chosen clue does not trigger any distractor words (mail, berlin, club, cap, web, train) before submitting.} \\
            }
        }
    }
    & & \\ \\

\end{supertabular}
}

\end{document}
